% SPDX-License-Identifier: Apache-2.0
\section{The Model of Time Enforces Itself}
\label{sec:e-time}

The merged specification spends a section on one rule. A timeless body may
not count cycles; a sequenced body counts them exactly as written. It gives
the two bodies two keywords, \code{flow} and \code{seq}, and states that a
cycle count inside a \code{flow} is a compile error.

Embedded in Rust, that rule needs no enforcement, because it cannot be
broken.

A cycle count is spelled \code{.await}. A plain \code{fn} is not
\code{async}, so it has no \code{.await} to write, so it cannot state a
cycle count. The compiler rejects the attempt already, with its own error
about awaiting outside an async context, and the specification says
nothing.

\begin{lstlisting}
// Timeless. There is no .await to write.
fn mul(a: u32, b: u32) -> u64 {
    a as u64 * b as u64
}

// Sequenced. Counts cycles exactly as written.
async fn main(&mut self) {
    loop {
        let req = self.host.request.recv().await;
        cycles(BAUD_TICKS).await;
    }
}
\end{lstlisting}

Two constructs are deleted, and a rule that needed a paragraph of
specification and a compiler check becomes a fact about the type system.

This is the clearest case of the reduction paying for itself, and it is
worth saying why it works rather than treating it as luck. The original
rule and Rust's rule are the same rule. LHDL's eighth thesis~\cite{theses}
argued that a designer should state intent and leave the mapping to time to
the implementation. Rust's \code{async} exists to mark exactly the place
where a computation gives up control and resumes later. Both are drawing
the same line, so one of them can be deleted.
